\subsection{Определения и основные свойства. Формула умножения. Преобразование Фурье свертки.}

\begin{definition}[Преобразование Фурье в \(L_1(\mathbb R)\)]
	Пусть
	\[
	f\in L_1(\mathbb R).
	\]
	Преобразованием Фурье функции \(f\) называется функция
	\[
	\widehat f(y)
	=
	\int_{\mathbb R} f(x)e^{-ixy}\,dx.
	\]
\end{definition}

\begin{theorem}[Основные свойства преобразования Фурье]
	Для любого \(f\in L_1(\mathbb R)\):
	\begin{enumerate}
		\item
		\[
		\|\widehat f\|_\infty\leq \|f\|_{L_1};
		\]
		
		\item \(\widehat f\) непрерывна;
		
		\item
		\[
		\widehat f(y)\to0,
		\qquad |y|\to\infty.
		\]
	\end{enumerate}
	
	Таким образом,
	\[
	\mathcal F:L_1(\mathbb R)\to C_0(\mathbb R)
	\]
	--- линейный ограниченный оператор.
\end{theorem}

\begin{proof}
	\begin{enumerate}
		\item
		\textbf{Хотим оценить норму \(\widehat f\).}
		
		Для любого \(y\in\mathbb R\)
		\[
		|\widehat f(y)|
		=
		\left|
		\int_{\mathbb R}f(x)e^{-ixy}\,dx
		\right|
		\leq
		\int_{\mathbb R}|f(x)|\,dx
		=
		\|f\|_{L_1}.
		\]
		Следовательно,
		\[
		\|\widehat f\|_\infty\leq\|f\|_{L_1}.
		\]
		
		\item
		\textbf{Хотим доказать непрерывность \(\widehat f\).}
		
		Пусть \(y_n\to y\). Тогда
		\[
		f(x)e^{-ixy_n}\to f(x)e^{-ixy}
		\]
		для почти всех \(x\), причем
		\[
		|f(x)e^{-ixy_n}|=|f(x)|\in L_1.
		\]
		По теореме Лебега
		\[
		\widehat f(y_n)\to\widehat f(y).
		\]
		
		\item
		\textbf{Хотим показать убывание на бесконечности.}
		
		Для \(\varphi\in C_c^1(\mathbb R)\) интегрированием по частям получаем
		\[
		\widehat\varphi(y)
		=
		\frac{1}{iy}
		\int_{\mathbb R}\varphi'(x)e^{-ixy}\,dx,
		\]
		поэтому
		\[
		|\widehat\varphi(y)|
		\leq
		\frac{\|\varphi'\|_{L_1}}{|y|}
		\to0.
		\]
		
		Теперь для \(f\in L_1\) выбираем
		\[
		\varphi\in C_c^1(\mathbb R),
		\qquad
		\|f-\varphi\|_{L_1}<\varepsilon.
		\]
		Тогда
		\[
		|\widehat f(y)|
		\leq
		|\widehat{f-\varphi}(y)|+|\widehat\varphi(y)|
		\leq
		\varepsilon+|\widehat\varphi(y)|.
		\]
		При достаточно больших \(|y|\)
		\[
		|\widehat\varphi(y)|<\varepsilon,
		\]
		следовательно,
		\[
		|\widehat f(y)|<2\varepsilon.
		\]
	\end{enumerate}
\end{proof}

\begin{corollary}
	Если
	\[
	f_n\to f
	\quad\text{в }L_1,
	\]
	то
	\[
	\widehat f_n\to\widehat f
	\quad\text{равномерно на }\mathbb R.
	\]
\end{corollary}

\begin{proof}
	Действительно,
	\[
	\|\widehat f_n-\widehat f\|_\infty
	=
	\|\widehat{f_n-f}\|_\infty
	\leq
	\|f_n-f\|_{L_1}
	\to0.
	\]
\end{proof}

\begin{theorem}[Формула умножения]
	Если
	\[
	f,g\in L_1(\mathbb R),
	\]
	то
	\[
	\int_{\mathbb R}\widehat f(y)g(y)\,dy
	=
	\int_{\mathbb R}f(x)\widehat g(x)\,dx.
	\]
\end{theorem}

\begin{proof}
	\textbf{Хотим поменять порядок интегрирования.}
	Для этого проверим абсолютную интегрируемость:
	\[
	\int_{\mathbb R}\int_{\mathbb R}
	|f(x)g(y)e^{-ixy}|\,dx\,dy
	=
	\|f\|_{L_1}\|g\|_{L_1}
	<\infty.
	\]
	По теореме Фубини
	\[
	\begin{aligned}
		\int_{\mathbb R}\widehat f(y)g(y)\,dy
		&=
		\int_{\mathbb R}
		\left(
		\int_{\mathbb R}f(x)e^{-ixy}\,dx
		\right)g(y)\,dy
		\\
		&=
		\int_{\mathbb R}f(x)
		\left(
		\int_{\mathbb R}g(y)e^{-ixy}\,dy
		\right)dx
		\\
		&=
		\int_{\mathbb R}f(x)\widehat g(x)\,dx.
	\end{aligned}
	\]
\end{proof}

\begin{definition}[Свертка]
	Пусть
	\[
	f,g\in L_1(\mathbb R).
	\]
	Их сверткой называется функция
	\[
	(f*g)(x)
	=
	\int_{\mathbb R}f(y)g(x-y)\,dy.
	\]
\end{definition}

\begin{theorem}
	Если
	\[
	f,g\in L_1(\mathbb R),
	\]
	то
	\[
	f*g\in L_1(\mathbb R)
	\]
	и
	\[
	\|f*g\|_{L_1}
	\leq
	\|f\|_{L_1}\|g\|_{L_1}.
	\]
\end{theorem}

\begin{proof}
	\textbf{Хотим показать интегрируемость свертки.}
	
	По теореме Тонелли
	\[
	\begin{aligned}
		\|f*g\|_{L_1}
		&\leq
		\int_{\mathbb R}\int_{\mathbb R}
		|f(y)||g(x-y)|\,dy\,dx
		\\
		&=
		\int_{\mathbb R}|f(y)|
		\left(
		\int_{\mathbb R}|g(x-y)|\,dx
		\right)dy
		\\
		&=
		\|f\|_{L_1}\|g\|_{L_1}.
	\end{aligned}
	\]
\end{proof}

\begin{theorem}[Преобразование Фурье свертки]
	Если
	\[
	f,g\in L_1(\mathbb R),
	\]
	то
	\[
	\widehat{f*g}
	=
	\widehat f\,\widehat g.
	\]
\end{theorem}

\begin{proof}
	\textbf{Хотим раскрыть свертку и поменять порядок интегрирования.}
	
	Так как
	\[
	\int_{\mathbb R}\int_{\mathbb R}
	|f(y)g(x-y)|\,dy\,dx
	=
	\|f\|_{L_1}\|g\|_{L_1}
	<\infty,
	\]
	по теореме Фубини
	\[
	\begin{aligned}
		\widehat{f*g}(z)
		&=
		\int_{\mathbb R}
		\left(
		\int_{\mathbb R}f(y)g(x-y)\,dy
		\right)e^{-ixz}\,dx
		\\
		&=
		\int_{\mathbb R}f(y)
		\left(
		\int_{\mathbb R}g(x-y)e^{-ixz}\,dx
		\right)dy.
	\end{aligned}
	\]
	
	\textbf{Хотим выделить преобразование Фурье функции \(g\).}
	Делаем замену
	\[
	t=x-y.
	\]
	Тогда
	\[
	\begin{aligned}
		\widehat{f*g}(z)
		&=
		\int_{\mathbb R}f(y)
		\left(
		\int_{\mathbb R}g(t)e^{-i(t+y)z}\,dt
		\right)dy
		\\
		&=
		\left(
		\int_{\mathbb R}f(y)e^{-iyz}\,dy
		\right)
		\left(
		\int_{\mathbb R}g(t)e^{-itz}\,dt
		\right)
		\\
		&=
		\widehat f(z)\widehat g(z).
	\end{aligned}
	\]
\end{proof}

\subsection{Преобразование Фурье в пространстве \texorpdfstring{\(L_2(\mathbb R)\)}{L2(R)}}

Рассмотрим пространство Шварца \(S(\mathbb R)\).
Используем без доказательства следующие факты:
\[
\overline{S(\mathbb R)}^{\,L_2}=L_2(\mathbb R),
\qquad
\mathcal F:S(\mathbb R)\to S(\mathbb R)
\]
--- биекция, причем
\[
\|\widehat f\|_{L_2}=\|f\|_{L_2},
\qquad f\in S(\mathbb R).
\]

Здесь
\[
\widehat f(y)
=
\frac{1}{\sqrt{2\pi}}
\int_{\mathbb R}f(x)e^{-ixy}\,dx.
\]

Так как \(S(\mathbb R)\) плотно в \(L_2(\mathbb R)\),
а \(\mathcal F\) является изометрией на \(S\),
она единственным образом продолжается до линейной изометрии
\[
\mathcal F:L_2(\mathbb R)\to L_2(\mathbb R).
\]

\begin{theorem}[Сохранение скалярного произведения]
	Для любых
	\[
	f,g\in L_2(\mathbb R)
	\]
	выполнено
	\[
	(f,g)=(\widehat f,\widehat g).
	\]
\end{theorem}

\begin{proof}
	\textbf{Сначала докажем равенство для \(f,g\in S\).}
	
	По формуле обращения
	\[
	\overline{g(x)}
	=
	\frac{1}{\sqrt{2\pi}}
	\int_{\mathbb R}
	\overline{\widehat g(y)}e^{-ixy}\,dy.
	\]
	
	По теореме Фубини
	\[
	\begin{aligned}
		(f,g)
		&=
		\int_{\mathbb R}f(x)\overline{g(x)}\,dx
		\\
		&=
		\frac{1}{\sqrt{2\pi}}
		\int_{\mathbb R}\int_{\mathbb R}
		f(x)\overline{\widehat g(y)}e^{-ixy}\,dy\,dx
		\\
		&=
		\int_{\mathbb R}
		\widehat f(y)\overline{\widehat g(y)}\,dy
		\\
		&=
		(\widehat f,\widehat g).
	\end{aligned}
	\]
	
	\textbf{Теперь переходим к \(f,g\in L_2\).}
	
	Берем
	\[
	f_n,g_n\in S,
	\qquad
	f_n\to f,
	\qquad
	g_n\to g
	\quad\text{в }L_2.
	\]
	
	Из непрерывности \(\mathcal F\)
	\[
	\widehat f_n\to\widehat f,
	\qquad
	\widehat g_n\to\widehat g.
	\]
	
	По непрерывности скалярного произведения
	\[
	(f_n,g_n)\to(f,g),
	\qquad
	(\widehat f_n,\widehat g_n)
	\to
	(\widehat f,\widehat g).
	\]
	
	Так как
	\[
	(f_n,g_n)
	=
	(\widehat f_n,\widehat g_n),
	\]
	переходом к пределу получаем
	\[
	(f,g)=(\widehat f,\widehat g).
	\]
\end{proof}

\begin{theorem}[Основные свойства преобразования Фурье в \(L_2\)]
	Для
	\[
	\mathcal F:L_2(\mathbb R)\to L_2(\mathbb R)
	\]
	выполнено:
	\begin{enumerate}
		\item
		\[
		\operatorname{Im}\mathcal F=L_2(\mathbb R);
		\]
		
		\item
		\[
		\ker\mathcal F=\{0\};
		\]
		
		\item \(\mathcal F\) не является компактным оператором;
		
		\item
		\[
		\mathcal F^4=I.
		\]
	\end{enumerate}
\end{theorem}

\begin{proof}
	\begin{enumerate}
		\item
		\textbf{Хотим доказать сюръективность.}
		
		Пусть
		\[
		f\in L_2.
		\]
		Так как \(S\) плотно в \(L_2\), существуют
		\[
		f_n\in S,
		\qquad
		f_n\to f.
		\]
		
		Так как
		\[
		\mathcal F:S\to S
		\]
		--- биекция, для каждого \(n\) существует
		\[
		g_n\in S:
		\qquad
		\widehat g_n=f_n.
		\]
		
		По изометричности
		\[
		\|g_n-g_m\|_{L_2}
		=
		\|\widehat g_n-\widehat g_m\|_{L_2}
		=
		\|f_n-f_m\|_{L_2}.
		\]
		
		Поэтому \(\{g_n\}\) фундаментальна, а значит
		\[
		g_n\to g
		\]
		для некоторого \(g\in L_2\).
		
		По непрерывности \(\mathcal F\)
		\[
		\widehat g_n\to\widehat g.
		\]
		Но
		\[
		\widehat g_n=f_n\to f,
		\]
		следовательно,
		\[
		\widehat g=f.
		\]
		
		Значит,
		\[
		\operatorname{Im}\mathcal F=L_2.
		\]
		
		\item
		\textbf{Хотим доказать инъективность.}
		
		Если
		\[
		\widehat f=0,
		\]
		то по изометричности
		\[
		\|f\|_{L_2}
		=
		\|\widehat f\|_{L_2}
		=
		0.
		\]
		Следовательно,
		\[
		f=0,
		\]
		то есть
		\[
		\ker\mathcal F=\{0\}.
		\]
		
		\item
		\textbf{Хотим доказать, что \(\mathcal F\) не компактен.}
		
		Предположим противное: \(\mathcal F\) компактен.
		Тогда композиция
		\[
		\mathcal F^4
		\]
		также компактна.
		
		Но
		\[
		\mathcal F^4=I.
		\]
		
		Получаем, что тождественный оператор
		\[
		I:L_2\to L_2
		\]
		компактен, что невозможно, поскольку \(L_2(\mathbb R)\)
		бесконечномерно.
		
		Следовательно, \(\mathcal F\) не компактен.
		
		\item
		\textbf{Хотим вычислить четвертую степень \(\mathcal F\).}
		
		На пространстве \(S\) из формулы обращения следует
		\[
		\mathcal F^2f(x)=f(-x).
		\]
		
		Поэтому
		\[
		\mathcal F^4f=f,
		\qquad f\in S.
		\]
		
		Так как \(S\) плотно в \(L_2\), а \(\mathcal F\) непрерывен,
		равенство продолжается на всё \(L_2\):
		\[
		\mathcal F^4=I.
		\]
	\end{enumerate}
\end{proof}

\begin{idea}
	\textbf{Зачем вообще нужно пространство Шварца и что мы здесь делаем.}
	
	В пространстве \(L_1(\mathbb R)\) преобразование Фурье можно определить напрямую:
	\[
	\widehat f(y)
	=
	\int_{\mathbb R}f(x)e^{-ixy}\,dx,
	\]
	поскольку
	\[
	\int_{\mathbb R}|f(x)e^{-ixy}|\,dx
	=
	\int_{\mathbb R}|f(x)|\,dx
	<\infty.
	\]
	
	\textbf{В \(L_2\) так сделать нельзя.}
	
	Из
	\[
	f\in L_2(\mathbb R)
	\]
	вообще не следует, что
	\[
	f\in L_1(\mathbb R).
	\]
	Поэтому для произвольного \(f\in L_2\) интеграл
	\[
	\int_{\mathbb R}f(x)e^{-ixy}\,dx
	\]
	может не существовать как обычный интеграл.
	
	То есть наша цель:
	\[
	\boxed{\text{определить преобразование Фурье для всех }f\in L_2}
	\]
	но напрямую по интегральной формуле это сделать нельзя.
	
	\medskip
	
	\textbf{1. Поэтому сначала берем ``хорошие'' функции.}
	
	Рассматриваем пространство Шварца
	\[
	S(\mathbb R)
	=
	\left\{
	f\in C^\infty(\mathbb R):
	\sup_{x\in\mathbb R}|x^m f^{(n)}(x)|<\infty
	\quad
	\forall m,n\geq0
	\right\}.
	\]
	
	Смысл условия: функции из \(S\) гладкие и вместе со всеми производными
	убывают на бесконечности быстрее любой степени.
	
	В частности,
	\[
	S(\mathbb R)\subset L_1(\mathbb R)\cap L_2(\mathbb R).
	\]
	
	Поэтому для \(f\in S\) интегральная формула преобразования Фурье
	точно имеет смысл:
	\[
	\widehat f(y)
	=
	\frac{1}{\sqrt{2\pi}}
	\int_{\mathbb R}f(x)e^{-ixy}\,dx.
	\]
	
	То есть сначала мы умеем корректно работать с
	\[
	\mathcal F:S\to S.
	\]
	
	\medskip
	
	\textbf{2. Почему именно \(S\), а не какое-то случайное подпространство?}
	
	У пространства Шварца есть два нужных нам свойства,
	которые используем без доказательства:
	\[
	\boxed{\overline{S}^{\,L_2}=L_2}
	\]
	и
	\[
	\boxed{\mathcal F:S\to S\text{ --- биекция}.}
	\]
	
	Первое означает, что любую функцию
	\[
	f\in L_2
	\]
	можно приблизить функциями из \(S\):
	\[
	f_n\in S,
	\qquad
	f_n\to f
	\quad\text{в }L_2.
	\]
	
	Поэтому, если мы хорошо определим преобразование Фурье на \(S\)
	и докажем его непрерывность в норме \(L_2\),
	то сможем распространить его на всё \(L_2\).
	
	\medskip
	
	\textbf{3. Главное свойство на \(S\) --- Фурье сохраняет норму.}
	
	Для \(f\in S\) доказывается равенство Планшереля:
	\[
	\|\widehat f\|_{L_2}
	=
	\|f\|_{L_2}.
	\]
	
	То есть
	\[
	\mathcal F:S\to S
	\]
	является изометрией.
	
	Это именно то, что позволяет перейти от \(S\) ко всему \(L_2\).
	
	\medskip
	
	\textbf{4. Как по \(S\) определить Фурье произвольной функции \(f\in L_2\)?}
	
	Берем последовательность хороших функций
	\[
	f_n\in S,
	\qquad
	f_n\to f
	\quad\text{в }L_2.
	\]
	
	У каждой \(f_n\) преобразование Фурье уже определено обычным интегралом:
	\[
	\widehat f_n
	=
	\mathcal Ff_n.
	\]
	
	Хотим определить
	\[
	\widehat f
	\]
	как предел этих преобразований.
	
	Из изометричности
	\[
	\|\widehat f_n-\widehat f_m\|_{L_2}
	=
	\|f_n-f_m\|_{L_2}.
	\]
	
	Так как
	\[
	f_n\to f,
	\]
	последовательность \(\{f_n\}\) фундаментальна.
	Следовательно, фундаментальна и
	\[
	\{\widehat f_n\}.
	\]
	
	Пространство \(L_2\) полно, поэтому существует
	\[
	g\in L_2:
	\qquad
	\widehat f_n\to g.
	\]
	
	Именно этот предел и объявляем преобразованием Фурье функции \(f\):
	\[
	\boxed{
		\mathcal Ff
		:=
		\lim_{n\to\infty}\mathcal Ff_n
	}.
	\]
	
	То есть для общей функции \(f\in L_2\)
	преобразование Фурье определяется уже
	\textbf{не обязательно обычным интегралом},
	а пределом в норме \(L_2\).
	
	\medskip
	
	\textbf{5. Почему результат не зависит от того, как именно мы приблизили \(f\)?}
	
	Пусть есть две последовательности
	\[
	f_n\to f,
	\qquad
	g_n\to f,
	\qquad
	f_n,g_n\in S.
	\]
	
	Тогда
	\[
	\|f_n-g_n\|_{L_2}\to0.
	\]
	
	По изометричности
	\[
	\|\widehat f_n-\widehat g_n\|_{L_2}
	=
	\|f_n-g_n\|_{L_2}
	\to0.
	\]
	
	Значит, обе последовательности преобразований дают один и тот же предел.
	
	Следовательно, \(\mathcal Ff\) определено корректно.
	
	\medskip
	
	\textbf{Итого вся конструкция выглядит так:}
	\[
	\boxed{
		S
		\overset{\text{плотно}}{\subset}
		L_2
	}
	\]
	\[
	\boxed{
		\mathcal F:S\to S
		\text{ уже определено обычным интегралом}
	}
	\]
	\[
	\boxed{
		\|\mathcal Ff\|_2=\|f\|_2
	}
	\]
	\[
	\Downarrow
	\]
	\[
	\boxed{
		\mathcal F
		\text{ единственным образом продолжается }
		S\to S
		\quad\text{до}\quad
		L_2\to L_2.
	}
	\]
	
	После этого мы уже изучаем свойства полученного оператора на \(L_2\):
	\[
	(f,g)=(\widehat f,\widehat g),
	\qquad
	\ker\mathcal F=\{0\},
	\qquad
	\operatorname{Im}\mathcal F=L_2,
	\qquad
	\mathcal F^4=I,
	\]
	и отсюда, в частности, следует, что \(\mathcal F\) не компактен.
	
	Таким образом, пространство Шварца здесь нужно не само по себе:
	это просто удобное плотное множество ``хороших'' функций,
	на котором преобразование Фурье сначала можно определить
	и изучить обычными интегралами, а затем по непрерывности
	перенести на всё \(L_2\).
\end{idea}